\documentclass[10pt,a4paper]{amsart}
\usepackage{amsmath,amssymb,amsthm,graphicx}
% For commutative diagrams we use Xy-pic
\usepackage[all]{xy}

% We use 2cell for 2-commutative diagrams.
\xyoption{2cell}
% Theorem environments.
%
\theoremstyle{plain}
\newtheorem{theorem}[subsection]{Theorem}
\newtheorem{proposition}[subsection]{Proposition}
\newtheorem{lemma}[subsection]{Lemma}

\theoremstyle{definition}
\newtheorem{definition}[subsection]{Definition}
\newtheorem{example}[subsection]{Example}
\newtheorem{exercise}[subsection]{Exercise}
\newtheorem{situation}[subsection]{Situation}

\theoremstyle{remark}
\newtheorem{remark}[subsection]{Remark}
\newtheorem{remarks}[subsection]{Remarks}

\numberwithin{equation}{subsection}

% Macros
%
\def\lim{\mathop{\mathrm{lim}}\nolimits}
\def\colim{\mathop{\mathrm{colim}}\nolimits}
\def\Spec{\mathop{\mathrm{Spec}}}
\def\Hom{\mathop{\mathrm{Hom}}\nolimits}
\def\Ext{\mathop{\mathrm{Ext}}\nolimits}
\def\SheafHom{\mathop{\mathcal{H}\!\mathit{om}}\nolimits}
\def\SheafExt{\mathop{\mathcal{E}\!\mathit{xt}}\nolimits}
\def\Sch{\mathit{Sch}}
\def\Mor{\mathop{\mathrm{Mor}}\nolimits}
\def\Ob{\mathop{\mathrm{Ob}}\nolimits}
\def\Sh{\mathop{\mathit{Sh}}\nolimits}
\def\NL{\mathop{N\!L}\nolimits}
\def\CH{\mathop{\mathrm{CH}}\nolimits}
\def\proetale{{pro\text{-}\acute{e}tale}}
\def\etale{{\acute{e}tale}}
\def\QCoh{\mathit{QCoh}}
\def\Ker{\mathop{\mathrm{Ker}}}
\def\Im{\mathop{\mathrm{Im}}}
\def\Coker{\mathop{\mathrm{Coker}}}
\def\Coim{\mathop{\mathrm{Coim}}}

% Boxtimes
%
\DeclareMathSymbol{\boxtimes}{\mathbin}{AMSa}{"02}

%
% Macros for moduli stacks/spaces
%
\def\QCohstack{\mathcal{QC}\!\mathit{oh}}
\def\Cohstack{\mathcal{C}\!\mathit{oh}}
\def\Spacesstack{\mathcal{S}\!\mathit{paces}}
\def\Quotfunctor{\mathrm{Quot}}
\def\Hilbfunctor{\mathrm{Hilb}}
\def\Curvesstack{\mathcal{C}\!\mathit{urves}}
\def\Polarizedstack{\mathcal{P}\!\mathit{olarized}}
\def\Complexesstack{\mathcal{C}\!\mathit{omplexes}}
% \Pic is the operator that assigns to X its picard group, usage \Pic(X)
% \Picardstack_{X/B} denotes the Picard stack of X over B
% \Picardfunctor_{X/B} denotes the Picard functor of X over B
\def\Pic{\mathop{\mathrm{Pic}}\nolimits}
\def\Picardstack{\mathcal{P}\!\mathit{ic}}
\def\Picardfunctor{\mathrm{Pic}}
\def\Deformationcategory{\mathcal{D}\!\mathit{ef}}
\usepackage{cleveref}
\crefname{originalsection}{original section}{original sections}
\Crefname{originalsection}{Original section}{Original sections}
\makeatletter
\newlabel{categories-example-2-fibre-product-categories}{{003R}{source}{}{sourcepointer-1}{}}
\newlabel{categories-example-2-fibre-product-categories@cref}{{[originalsection][0][]003R}{[1][0][]source}{}{}{}}
\makeatother
\crefname{originalreference}{original reference}{original references}
\Crefname{originalreference}{Original reference}{Original references}
\makeatletter
\newlabel{lemma-proper-base-change-f-star}{{0A3U}{source}{}{sourcepointer-2}{}}
\newlabel{lemma-proper-base-change-f-star@cref}{{[originalreference][0][]0A3U}{[1][0][]source}{}{}{}}
\makeatother
\title{A proper modification that is an isomorphism outside a closed subscheme gives a fibre product decomposition of etale sheaves}
\author{The Stacks Project Authors}
\date{stacks-project:etale-cohomology.tex:local-snapshot-20261011; GNU FDL 1.2 or later}
\begin{document}
\maketitle
\noindent\textit{Editorial note.} All delivered mathematical text is the author's own, copied from the pinned Stacks Project chapter file etale-cohomology.tex; no statement or proof step was rewritten, reordered or removed.\par\medskip
\noindent\textit{Editorial note.} The counted proof is the author's complete argument: the proof is the author's construction and stalk computation in full; its closing sentence is quoted verbatim here: "To finish the proof, we observe that there are canonical maps i\textasciicircum{}\{-1\}F -> G and f\textasciicircum{}\{-1\}F -> F' compatible with alpha which on stalks produce the isomorphisms we saw above. We omit the careful construction of these maps." That omitted construction is the routine bookkeeping of the final verification, all substance (the quasi-inverse functor and the two stalk computations) being written out.\par\medskip
\noindent\textit{Editorial note.} The label lemma-glue-etale-sheaf-modification names the counted statement, and the proof environment immediately following it in the source is delivered verbatim as its proof.\par\medskip
\noindent\textit{Editorial note.} The counted statement is carried by the author's own lemma environment in the chapter source.\par\medskip
\noindent\textit{Editorial note.} The statement's cartesian diagram is the author's and is delivered verbatim; the two cited inputs are kept as external pointers with their canonical Stacks tags.\par\medskip
\noindent\textit{Editorial note.} Pre-open gate (§17.1): the label was checked against the lane list, all focused-seeds claim\_keys, and the published items' claim\_keys before the card was written.\par\medskip
\noindent\textit{Allowed standard prerequisites (editorial).}\begin{itemize}
\item The 2-fibre product of categories. Categories, the example constructing the 2-fibre product category used to state the target of the functor. Supplies the category of triples in which the equivalence is asserted.
\item Proper base change for the pushforward. Same chapter, the lemma identifying i\textasciicircum{}\{-1\}f\_*F' with g\_*j\textasciicircum{}\{-1\}F' for the cartesian square of the statement. Supplies the isomorphism used to build the quasi-inverse functor and again in the stalk computation over Z.
\end{itemize}\medskip
\ifdefined\hypertarget\hypertarget{sourcepointer-1}{}\fi
\noindent\textit{Source reference.} Original section 003R: sources/stacks-chow/tags.txt, source lines 145--145. Prerequisite pointer invoked at proof line(s) 23187; it names a standard fact of a sibling chapter that the counted argument uses, and every step of the counted proof built on it is retained verbatim in the delivered body.\par\medskip
\ifdefined\hypertarget\hypertarget{sourcepointer-2}{}\fi
\noindent\textit{Source reference.} Original reference 0A3U: sources/stacks-etale-cohomology/pinned-original-source.tex, source lines 18282--18282. Prerequisite pointer invoked at proof line(s) 23200; it names a standard fact of the same chapter that the counted argument uses, and every step of the counted proof built on it is retained verbatim in the delivered body.\par\medskip
\begin{lemma}
\label{lemma-glue-etale-sheaf-modification}
Let $f : X' \to X$ be a proper morphism of schemes. Let $i : Z \to X$
be a closed immersion. Set $E = Z \times_X X'$. Picture
$$
\xymatrix{
E \ar[d]_g \ar[r]_j & X' \ar[d]^f \\
Z \ar[r]^i & X
}
$$
If $f$ is an isomorphism over $X \setminus Z$, then the functor
$$
\Sh(X_\etale)
\longrightarrow
\Sh(X'_\etale) \times_{\Sh(E_\etale)} \Sh(Z_\etale)
$$
is an equivalence of categories.
\end{lemma}

\begin{proof}
We will work with the $2$-fibre product category as constructed in
Categories, Example \ref{categories-example-2-fibre-product-categories}.
The functor sends $\mathcal{F}$ to the triple
$(f^{-1}\mathcal{F}, i^{-1}\mathcal{F}, c)$ where
$c : j^{-1}f^{-1}\mathcal{F} \to g^{-1}i^{-1}\mathcal{F}$
is the canonical isomorphism. We will construct a quasi-inverse functor. Let
$(\mathcal{F}', \mathcal{G}, \alpha)$ be an object
of the right hand side of the arrow.
We obtain an isomorphism
$$
i^{-1}f_*\mathcal{F}' = g_*j^{-1}\mathcal{F}'
\xrightarrow{g_*\alpha}
g_*g^{-1}\mathcal{G}
$$
The first equality is Lemma \ref{lemma-proper-base-change-f-star}.
Using this we obtain maps
$i_*\mathcal{G} \to i_*g_*g^{-1}\mathcal{G}$
and $f'_*\mathcal{F}' \to i_*g_*g^{-1}\mathcal{G}$. We set
$$
\mathcal{F} = f_*\mathcal{F}' \times_{i_*g_*g^{-1}\mathcal{G}} i_*\mathcal{G}
$$
and we claim that $\mathcal{F}$ is an object of the left hand side
of the arrow whose image in the right hand side is isomorphic to
the triple we started out with. Let us compute the stalk of $\mathcal{F}$
at a geometric point $\overline{x}$ of $X$.

\medskip\noindent
If $\overline{x}$ is not
in $Z$, then on the one hand $\overline{x}$ comes from a unique
geometric point $\overline{x}'$ of $X'$ and
$\mathcal{F}'_{\overline{x}'} = (f_*\mathcal{F}')_{\overline{x}}$
and on the other hand we have $(i_*\mathcal{G})_{\overline{x}}$
and $(i_*g_*g^{-1}\mathcal{G})_{\overline{x}}$ are singletons.
Hence we see that $\mathcal{F}_{\overline{x}}$ equals
$\mathcal{F}'_{\overline{x}'}$.

\medskip\noindent
If $\overline{x}$ is in $Z$, i.e., $\overline{x}$ is the image of
a geometric point $\overline{z}$ of $Z$, then we obtain
$(i_*\mathcal{G})_{\overline{x}} = \mathcal{G}_{\overline{z}}$
and
$$
(i_*g_*g^{-1}\mathcal{G})_{\overline{x}} =
(g_*g^{-1}\mathcal{G})_{\overline{z}} =
\Gamma(E_{\overline{z}}, g^{-1}\mathcal{G}|_{E_{\overline{z}}})
$$
(by the proper base change for pushforward used above)
and similarly
$$
(f_*\mathcal{F}')_{\overline{x}} =
\Gamma(X'_{\overline{x}}, \mathcal{F}'|_{X'_{\overline{x}}})
$$
Since we have the identification
$E_{\overline{z}} = X'_{\overline{x}}$ and since $\alpha$
defines an isomorphism between the sheaves
$\mathcal{F}'|_{X'_{\overline{x}}}$ and
$g^{-1}\mathcal{G}|_{E_{\overline{z}}}$
we conclude that we get
$$
\mathcal{F}_{\overline{x}} = \mathcal{G}_{\overline{z}}
$$
in this case.

\medskip\noindent
To finish the proof, we observe that there are canonical maps
$i^{-1}\mathcal{F} \to \mathcal{G}$ and $f^{-1}\mathcal{F} \to \mathcal{F}'$
compatible with $\alpha$ which on stalks produce the isomorphisms
we saw above. We omit the careful construction of these maps.
\end{proof}

\end{document}
